\Needspace{15\baselineskip}
\fancyhead[L]{\small JUNO ELEMENTE VIII.}
\begin{center}
JUNO ELEMENTE VIII.
\end{center}

\hrule
\begin{center}
\fontsize{10}{12}\selectfont Göttingische gelehrte Anzeigen. Stück 136. Seite 1353..1355. 25. August 1808.
\end{center}
\par\vspace{4pt}\hrule\vspace{4pt}

\par\indent Die ersten Beobachtungen der \emph{Juno}, die in diesem Jahr an der hiesigen Sternwarte gemacht wurden, sind bereits im 107. Stück dieser Zeitschrift [S. 303 dieses Bandes] mitgeteilt worden. Seitdem wurden die Beobachtungen fortgesetzt, sooft die Umstände günstig waren, freilich nur mit dem Kreismikrometer, da der kleine Planet im lichtschwachen Fernrohr unseres Mauerquadranten durchweg unsichtbar blieb. Man hat versucht, diesen Mangel durch zahlreiche und möglichst sorgfältige Vergleiche auszugleichen. Besonders die vier Beobachtungen in der Nähe der Opposition dürfen für so genau gehalten werden, wie es das angewandte Verfahren überhaupt zulässt. Da die Beobachtung vom 22. Juni noch eine kleine Korrektur erfahren hat, stellen wir hier die ersten beiden Bestimmungen noch einmal mit allen späteren zusammen, die bisher gemacht wurden [sie stehen oben auf Seite 305 im Brief an \textsc{von Lindenau} vom 6. Aug. 1808].

\par\indent Demnach ist der Fehler der Ephemeride bei den letzten Beobachtungen in der Rektaszension auf 13 Minuten angewachsen. Der Fehler in der Deklination liegt ziemlich unverändert bei $1\frac{1}{2}$ Minuten. Beobachtungen von anderen Orten sind bisher noch nicht bekannt geworden.

\par\indent Professor \textsc{Gauss} hat die letzten vier Beobachtungen, die besonders gut ausgefallen sind, zur Bestimmung der Opposition benutzt und dabei folgendes Ergebnis erhalten:

\begin{center}
1808. August 2. $9^{\text{u}}\,30^{\text{m}}\,43^{\text{s}}$ mittlere Zeit in Göttingen\\
wahre Länge .\;.\;.\;.\;.\;.\;.\;.\;.\; $310^\circ\,16'\,31''5$\\
wahre geozentrische Breite .\;.\;.\;.\; $13\;53\;56,2$ nordl.
\end{center}

\Needspace{25\baselineskip}
\par\indent Diese Opposition, die \emph{vierte} bisher beobachtete, wurde anschließend mit denen von 1804, 1806 und 1807 verbunden, um die folgenden neuen Elemente (VIII) zu bestimmen. Sie geben sämtliche bisherigen Beobachtungen weiterhin sehr gut wieder.


\begin{center}\fontsize{10}{12}\selectfont\setlength{\tabcolsep}{3pt}\renewcommand{\arraystretch}{1.15}\begin{tabular}{lr}
\multicolumn{2}{c}{Epoche der mittleren Länge}\\
\multicolumn{2}{c}{für den Meridian von Göttingen:}\\\hline
1804 .\;.\;.\;.\;.\;.\;.\;.\; & $320^\circ\;1'\;20''1$\\
1805 .\;.\;.\;.\;.\;.\;.\;.\; & $42\;35\;8.4$\\
1806 .\;.\;.\;.\;.\;.\;.\;.\; & $125\;8\;56.7$\\
1807 .\;.\;.\;.\;.\;.\;.\;.\; & $207\;42\;45.0$\\
1808 .\;.\;.\;.\;.\;.\;.\;.\; & $290\;30\;7.6$\\
1809 .\;.\;.\;.\;.\;.\;.\;.\; & $13\;3\;55.8$\\
1810 .\;.\;.\;.\;.\;.\;.\;.\; & $95\;37\;44.1$\\
\end{tabular}\end{center}

\begin{center}\fontsize{9}{11}\selectfont\setlength{\tabcolsep}{3pt}\renewcommand{\arraystretch}{1.15}\begin{tabular}{lr}
Tägliche mittlere tropische Bewegung .\;.\; & $814''324$\\
Tropische Umlaufszeit .\;.\; & 1591 Tage 12 Stunden\\
Perihel 1805 .\;.\;.\;.\;.\;.\;.\; & $53^\circ\;10'\;53''9$\\
Aufsteigender Knoten 1805 .\;.\;.\;.\; & $171\;4\;11{,}3$\\
\multicolumn{2}{l}{\quad Beide als fest gegenüber den Fixsternen vorausgesetzt.}\\
Neigung der Bahn .\;.\;.\;.\;.\;.\;.\;.\; & $13\;4\;11{,}0$\\
Exzentrizität .\;.\;.\;.\;.\;.\;.\;.\;.\;.\;.\; & $0{,}2554521$\\
Logarithmus der großen Halbachse .\;.\;.\; & $0{,}4261883$\\
\end{tabular}\end{center}


\SourceLeaf{312}\par\indent Nach diesen Elementen fällt die nächste Opposition, bei der Juno wieder beträchtlich heller sein wird, auf den 30. Januar 1810, bei $130^\circ\,2'$ Länge und $14^\circ\,51'$ südlicher Breite am Kopf der Wasserschlange.

